\section{Ordered renewal probes}
\label{sec:10-thermal-probes}

Two adjacent pieces of a self-avoiding path cannot pass through each
other. We quantify this elementary obstruction under the independent
renewal law before imposing the eventual height or length of either
piece. The estimate must also allow a single long irreducible bridge to
be inspected from both ends. Those two inspections share the same
critical weight, so their costs cannot be multiplied as independent
probabilities.

Throughout this section, heights count lattice bands, port paths carry
weight $\rho^{|\gamma|}$, and $|\gamma|$ counts visited vertices. Put
$\alpha=3/4$. We use the bridge masses
$B_j\asymp(1+j)^{\alpha-1}$, with $B_0=1$, the localization and kernel
variation estimates of Lemmas~\ref{lem:08-localization}
and~\ref{lem:08-variation}, the free-end tubes of
Lemma~\ref{lem:08-free-tube}, and the corridor construction of
Corollary~\ref{lem:08-long-arch}. A free endpoint is summed over its lattice
ports; once its shape is specified, attaching it to another path makes
its translation unique.

The argument has three parts. First we fix the probability and
unnormalized measures used in every probe. Next a recoverable hairpin
construction bounds masses of pairs with suitably separated renewal
endpoints. Finally we show that, except for a controlled exceptional
part, an ordered probe supplies many such endpoints. This last step
uses its conditional failure bound before conditioning on continued
avoidance.

\subsection{Stopped paths and renewal occurrences}

Write \(\nu(I)=\rho^{|I|}\) for the probability mass of a rooted
irreducible strict port bridge, with the endpoint convention already used
in bridge factorization. Thus \(\sum_I\nu(I)=1\). Let \(\mathbf P\) be
the product law of an infinite sequence of such bridges, with renewal
heights \(U_0=0,U_1,\ldots\). For every real threshold $t\ge0$, let \(\mathcal Q_t\) be the law of its
prefix through the first renewal at or above \(t\). For \(t=0\) this is
the empty prefix. In contrast, define the renewal-occurrence measure
\(\mathcal R\) on finite strings by
\[
 \int G(P)\,d\mathcal R(P)
   =\mathbf E\sum_{k\ge0}G(I_1\cdots I_k).
\]
At fixed height \(j\) this is exactly the unnormalized bridge measure;
its endpoint marginal is \(b_j(x)\), summed with counting measure over
lattice ports \(x\), and its total mass is \(B_j\).
In particular \(\mathcal Q_t\) is a probability measure, whereas
\(\mathcal R\{P:\operatorname{ht}(P)\le t\}\asymp(1+t)^\alpha\).

Whenever a renewal prefix of height \(v\ge t\) is marked, it can be cut
at its first renewal \(U\ge t\). Its exact measure is
\(d\mathcal Q_t(P)\,d\mathcal R(R)\), with the residual-height condition
\(U+\operatorname{ht}(R)=v\). For a sum over \(v\le Ct\), the residual
mass is uniformly at most \(C(1+t)^\alpha\). The empty residual string
is included. This identity, rather than conditioning a bridge to hit a
specified terminal height, is the change of measure used below.

\begin{lemma}[Jump charges and renewal counts]
\label{lem:10-charges}
For an irreducible bridge let \(J\) be its height and \(W\) its maximum
absolute horizontal displacement from its start. For dyadic \(a\ge1\) and real $d\ge0$,
\[
 \mathbf P\{a\le J<2a,\ W\ge da\}
 \le Ca^{-\alpha}\Phi(d),\qquad \Phi(d)=(1+d)^{-5/4}.
\]
Moreover, for $b\ge0$ and \(x\ge1\),
\[
 \mathbf E\sum_{k:U_{k-1}<b}
 \min\{1,(J_k+W_k)/x\}
 \le C(1+b)^\alpha x^{-\alpha}.
\]
For every positive integer $p$, the number of renewals in any height interval of length \(v\ge0\), after a stopping
renewal, has \(p\)-th moment at most
\(C_p(1+v)^{p\alpha}\), uniformly in the stopping history and interval
location.
\end{lemma}

\begin{proof}
Prepend and append bridges with heights in \([a,2a]\) to a marked
irreducible. Their product mass is at least \(ca^{2\alpha}\). Output
heights are \(O(a)\), and a marked irreducible of height at least \(a\)
has bounded multiplicity in such an output. Its lateral excursion forces
comparable output range. The strip lateral estimate bounds output mass
by \(Ca^\alpha\Phi(d)\), proving the first assertion. Summing dyadic
\(a\) and width bins gives
\(\mathbf E\min\{1,(J+W)/x\}\le Cx^{-\alpha}\).
The event \(U_{k-1}<b\) is decided before jump \(k\); multiplying by the
fresh-jump expectation and summing gives the charge estimate, because
\(\sum_{j<b}B_j\le C(1+b)^\alpha\). This includes the last, possibly
overshooting jump. For the moment estimate, order the marked renewals,
restart at each, and sum their nonnegative height increments using
\(\sum_{j\le v}B_j\le C(1+v)^\alpha\). Diagonal repetitions cost only
a constant depending on \(p\). A delayed interval is handled by first
stopping on entrance; overshoot can only shorten its remaining length.
\end{proof}

Launch two independent renewal strings at neighboring ports on one line.
The left and right designations specify their launch order. The test
\(E(h)\), measurable under \(\mathcal Q_h\otimes\mathcal Q_h\), requires
disjointness at depths \(0\le y<h\) and, at every such depth,
\[
 \begin{aligned}
 \max x(\text{left arc at }y)&\le\max x(\text{right arc at }y),\\
 \min x(\text{left arc at }y)&\le\min x(\text{right arc at }y).
 \end{aligned}
\]
It is enough to use depths intersecting the embedded edges. Strict
renewals prevent later visits below their level, so \(E(h)\) implies
all smaller-depth tests. This is an ordered test; disjointness alone of
the two stopped strings is not its definition.

Fix constants \(L_2\ge L>1\), to be selected below. A dyadic height is $2^k$ with $k\in\mathbb Z_{\ge0}$. For dyadic \(r\ge h\),
let \(F(h,r,d)\) be the probability of \(E(h)\) together with a jump
\(I\) of a designated arm satisfying \(r\le J(I)<2r\), \(W(I)\ge dr\),
and born before height \(h/L\). Such a jump is unique: after it the renewal height is at least $h$, so a second jump cannot be born in the specified window. Equivalently,
\begin{equation}
 F(h,r,d)=
 \int_{\operatorname{ht}(P)<h/L}\!d\mathcal R(P)
 \sum_{\substack{r\le J(I)<2r\\W(I)\ge dr}}\nu(I)
 \int \mathbf1_{E(h)}(PI,A)\,d\mathcal Q_h(A).
 \label{eq:10-F-integral}
\end{equation}
The obstacle \(PI\) already contains its stop at \(h\).

For the paired quantity use the same irreducible \(I\) from its two ends,
reversing and reflecting it as needed. Precede it, in these separate
inward orientations, by prefixes \(P_i\) of height \(f_i<h_i/L_2\),
and launch independent workers \(A_i\) at adjacent ports. There is no
avoidance condition between the two sides. Its precise definition is
\begin{equation}
 \begin{split}
 F_{12}(h_1,h_2,r)
 ={}&\sum_{r\le J(I)<2r}\nu(I)
 \int_{f_1<h_1/L_2}\!d\mathcal R(P_1)
 \int_{f_2<h_2/L_2}\!d\mathcal R(P_2)\\[-2mm]
 &\quad\times\prod_{i=1}^2
 \left[\int\mathbf1_{E(h_i)}(P_iI^{(i)},A_i)
                    d\mathcal Q_{h_i}(A_i)\right].
 \end{split}
 \label{eq:10-paired-integral}
\end{equation}
Here \(I^{(1)},I^{(2)}\) are the two normalized orientations of one shape,
not independent jumps. Each side is normalized by its own near launch;
its other ports are determined by concatenation. In particular no
horizontal position is summed in addition to the endpoints already
present in \(d\mathcal R\) and \(\nu\).

Figure~\ref{fig:10-probe-factors} displays the shared and separate
integrations in this definition.

\begin{figure}[ht]
\centering
\begin{tikzpicture}[font=\small,>=Stealth,
  mass/.style={draw,rounded corners,inner sep=5pt,fill=black!4},
  test/.style={draw,inner sep=4pt}]
 \node[mass] (giant) at (0,1.7) {$\nu(I)$, counted once};
 \node[test] (e1) at (-2,.55) {$E(h_1)$};
 \node[test] (e2) at (2,.55) {$E(h_2)$};
 \draw[->] (giant)--node[left] {$I^{(1)}$}(e1);
 \draw[->] (giant)--node[right] {$I^{(2)}$}(e2);
 \node[align=center] (p1) at (-2,-.65)
   {$P_1:\ f_1<h_1/L_2$\\$d\mathcal R(P_1)\,d\mathcal Q_{h_1}(A_1)$};
 \node[align=center] (p2) at (2,-.65)
   {$P_2:\ f_2<h_2/L_2$\\$d\mathcal R(P_2)\,d\mathcal Q_{h_2}(A_2)$};
 \draw[->] (p1)--(e1);
 \draw[->] (p2)--(e2);
\end{tikzpicture}
\caption{The integrations defining $F_{12}$. Both tests use one
irreducible bridge. Conditional on that shape, their prefix and worker
integrations separate. Prefix endpoint positions are already summed
by $\mathcal R$; there is no additional translation factor.}
\label{fig:10-probe-factors}
\end{figure}

\begin{theorem}[Ordered probes, with one or two giant ends]
\label{thm:10-ordered-probes}
The constants \(L,L_2,K_1,K_2,K_3\) can be chosen so that, for all integer
\(1\le h,h_1,h_2\le r\), with \(r\) dyadic and \(d\ge0\),
\begin{equation}
 \mathbf P(E(h))\le K_1h^{-\alpha},\qquad
 F(h,r,d)\le K_2r^{-\alpha}\Phi(d),\qquad
 F_{12}(h_1,h_2,r)\le K_3r^{-\alpha}.\label{eq:T-34}
\end{equation}
Either adjacent-port choice and each reflected orientation obey the same
bounds.
\end{theorem}

\subsection{Recovering one or two hairpins}

Fix \(e>0\), \(D<\infty\), and \(0<\sigma<e/10\). A worker bridge
\(A\), launched to the right of an obstacle, is \emph{qualified at scale
\(h\)} if its endpoint height \(u\) lies in \([eh,2h/5]\), the two
shapes are disjoint, \(|x|\le Dh\) on \(A\), and the obstacle satisfies
\[
 x\le Dh\quad(0\le y\le u+\sigma h),\qquad
 x\le x_A-\sigma h\quad(u\le y\le u+\sigma h).
\]
Coordinates are relative to the near launch, and \(x_A\) is the worker
endpoint abscissa. The obstacle may extend arbitrarily far left.
Reflected qualifications are allowed throughout.

Write $\chi_h(O,A)$ for the indicator that the worker bridge $A$
is qualified against the obstacle $O$, with the prescribed neighboring
launches. The three unnormalized masses used below are
\begin{align*}
 Y(h)&=\iint \chi_h(O,A)
 \mathbf1_{\{\operatorname{ht}(A)+2\sigma h\le
                    \operatorname{ht}(O)\le h\}}
 \,d\mathcal R(O)\,d\mathcal R(A),\\
 X_d(h,r)&=\sum_{\substack{r\le J(I)<2r\\W(I)\ge dr}}\nu(I)
 \int_{\operatorname{ht}(P)\le h/8}d\mathcal R(P)
 \int\chi_h(PI,A)\,d\mathcal R(A),\\
 X_{12}(h_1,h_2,r)&=
 \sum_{r\le J(I)<2r}\nu(I)
 \prod_{i=1}^2\left[
 \int_{\operatorname{ht}(P_i)\le h_i/8}d\mathcal R(P_i)
 \int\chi_{h_i}(P_i I^{(i)},A_i)\,d\mathcal R(A_i)
 \right].
\end{align*}
The last line uses the two inward orientations of one irreducible,
exactly as in \eqref{eq:10-paired-integral}. Its factor $\nu(I)$
occurs once. Worker endpoints are summed by $\mathcal R$; these
quantities are masses of endpoint occurrences, not probabilities of
stopped paths. There is no extra horizontal-translation sum.

\begin{lemma}[Hairpin mass bounds]
\label{lem:10-hairpins}
For fixed qualification parameters and all sufficiently large scales,
with \(r\) dyadic, \(1\le h,h_1,h_2\le r\), and \(d\ge0\),
\[
 Y(h)\le Ch^\alpha,\qquad
 X_d(h,r)\le Ch^\alpha r^{-\alpha}\Phi(d),\qquad
 X_{12}(h_1,h_2,r)\le Ch_1^\alpha h_2^\alpha r^{-\alpha}.
\]
Bounded scales can be included by changing constants.
\end{lemma}

\begin{proof}
\emph{Constructing the hairpin.}
Apply Corollary~\ref{lem:08-long-arch} at the worker endpoint. It
supplies an arch contained in the next $\sigma h$ bands, to the right
of $x_A-\sigma h$, whose far endpoint lies beyond $D'h$, for any
sufficiently large fixed $D'$. To pass from the relative displacement
in that lemma to this absolute bound, increase its fixed displacement
constant by $D$, using $|x_A|\le Dh$. Its total critical mass is between
$c h^{-1/4}$ and $C h^{-1/4}$. The lower bound has no exponent loss;
the fixed-endpoint tube bound would not suffice here. The corridor
lemma also establishes the endpoint half-plane bounds needed when
this arch is cut in the second-moment calculation below.

Choose \(z\in[eh/5,eh/3]\). From the far end of the long arch descend by
a narrow bridge to level \(z\), with mass at least \(ch^{-1/4}\).
Continue downward by a bridge of free height in \([2r,4r]\); use \(r=h\)
for \(Y\). Require its portion through depth \(z+1\) from its starting
port to stay within horizontal distance \(D''h\) of that port. This
restriction retains mass at least \(cr^\alpha\). Indeed its failure
probability in the renewal law is small for large \(D''\), by
Lemma~\ref{lem:10-charges}. If \(N\) counts renewals in the desired
terminal interval, then \(\mathbf EN\ge cr^\alpha\) and
\(\mathbf EN^2\le Cr^{2\alpha}\); consequently
\(\mathbf E[N\mathbf1_{\rm failure}]\le
 C r^\alpha\mathbf P(\mathrm{failure})^{1/2}\).
Take \(D''\) large and then \(D'\) much larger.

Join the two near launches by a fixed short arch in the one band on the
opposite side. The resulting path rises from its new lower end, goes
round the long arch, descends along the worker to this short turn, then
rises along the obstacle. Append an upper continuation with height in
\([2r,4r]\). For \(X_{12}\) also make the reflected hairpin at the other
end. Qualification gives all lateral clearances. The two near-end
structures protrude by at most \((2/5+\sigma)h_i\); since \(h_i\le r\)
and the giant has height at least \(r\), they are separated in height.
These operations produce strict bridges at a fixed translated root.
Local connections change weights by fixed factors.
Figure~\ref{fig:10-hairpin-geometry} shows the one-giant construction;
there \(\tau\) is the absolute height of the near launch.
\begin{figure}[htbp]
\centering
\begin{tikzpicture}[x=1cm,y=1cm,font=\small,>=Stealth,
  pathline/.style={line width=.9pt,line cap=round,line join=round},
  guide/.style={gray!65,densely dashed,line width=.4pt}]
  \fill[gray!9] (.65,3.40) rectangle (8.20,3.88);
  \draw[guide] (.65,3.40)--(8.20,3.40);
  \draw[guide] (.65,3.88)--(8.20,3.88);
  \draw[guide] (3.15,3.40)--(3.15,3.88);
  \node[anchor=south,font=\footnotesize] at (3.15,3.94) {$x_A-\sigma h$};
  \draw[guide] (.65,1.22)--(8.20,1.22);
  \draw[guide] (5.95,1.72)--(8.20,1.72);
  \node[anchor=west] at (8.25,1.22) {$\tau$};
  \node[anchor=west] at (8.25,1.72) {$\tau+z$};
  \node[anchor=west] at (8.25,3.40) {$\tau+u$};
  \node[anchor=west] at (8.25,3.88) {$\tau+u+\sigma h$};
  \node[anchor=west,font=\footnotesize] at (8.25,4.38) {absolute heights};
  \draw[pathline] (2.35,1.22) .. controls (2.05,1.39) and (2.45,1.52) .. (2.12,1.72);
  \draw[pathline] (2.12,1.72)
    .. controls (1.48,1.82) and (1.02,2.28) .. (1.60,2.38)
    .. controls (2.43,2.55) and (2.26,1.96) .. (2.66,2.20)
    .. controls (3.08,2.64) and (1.20,2.80) .. (1.33,3.20)
    .. controls (1.48,3.68) and (2.51,3.11) .. (2.58,3.60)
    .. controls (2.68,4.13) and (1.07,3.99) .. (1.25,4.42)
    .. controls (1.51,4.96) and (4.95,4.23) .. (5.06,4.88)
    .. controls (5.15,5.30) and (3.80,5.40) .. (3.33,5.12)
    .. controls (2.55,4.69) and (2.27,5.25) .. (2.65,5.61);
  \draw[pathline] (2.65,5.61) .. controls (2.97,5.80) and (2.39,5.98) .. (2.70,6.30);
  \draw[->,pathline] (2.58,5.52)--(2.65,5.61);
  \draw[->,pathline] (2.56,6.10)--(2.70,6.30);
  \node[anchor=east] at (2.02,1.43) {$P$};
  \node[fill=white,inner sep=2pt] at (1.20,2.82) {$I$};
  \fill (2.12,1.72) circle (1.4pt);
  \fill (2.65,5.61) circle (1.4pt);
  \node[anchor=west] at (2.95,6.07) {upper exterior};
  \draw[pathline,blue!65!black] (2.82,1.22)
    .. controls (3.15,1.54) and (2.88,1.79) .. (3.18,2.04)
    .. controls (3.47,2.26) and (3.09,2.46) .. (3.39,2.68)
    .. controls (3.69,2.90) and (3.49,3.13) .. (3.68,3.40);
  \draw[->,pathline,blue!65!black] (3.58,3.24)--(3.51,3.08);
  \node[anchor=west,text=blue!65!black] at (3.57,2.57) {$\overline A$};
  \fill (3.68,3.40) circle (1.4pt);
  \node[anchor=north west] at (3.85,3.28) {worker endpoint};
  \draw[pathline] (2.82,1.22) .. controls (2.83,.85) and (2.32,.85) .. (2.35,1.22);
  \draw[->,pathline] (2.64,.96)--(2.50,.96);
  \node[anchor=north,align=center] at (2.58,.77) {fixed short turn\\below $\tau$};
  \draw[pathline,orange!75!black] (7.20,1.72)
    .. controls (7.39,2.14) and (7.05,2.75) .. (7.20,3.40);
  \draw[pathline,orange!75!black] (7.20,3.40)
    .. controls (7.17,3.76) and (6.19,3.78) .. (5.50,3.69)
    .. controls (4.81,3.61) and (3.82,3.86) .. (3.68,3.40);
  \draw[->,pathline,orange!75!black] (5.63,3.70)--(5.43,3.68);
  \draw[->,pathline,orange!75!black] (7.19,2.84)--(7.17,3.02);
  \node[anchor=south,text=orange!75!black] at (5.61,3.91) {long arch};
  \node[anchor=west,align=left,text=orange!75!black] at (7.47,2.73) {narrow\\bridge};
  \draw[pathline] (7.45,.13)
    .. controls (7.64,.43) and (6.86,.61) .. (7.14,.93)
    .. controls (7.42,1.24) and (6.92,1.42) .. (7.20,1.72);
  \draw[->,pathline] (7.10,1.57)--(7.20,1.72);
  \fill (7.20,1.72) circle (1.4pt);
  \node[anchor=west,align=left] at (7.70,.48) {lower\\exterior};
  \node[font=\footnotesize] at (5.28,1.57) {$z<u$};
  \node[anchor=west,font=\footnotesize] at (.63,6.64)
    {One-giant output: arrows show its traversal};
\end{tikzpicture}
\caption{The one-giant hairpin, schematically. From its new lower end the
output follows the right exterior to height $\tau+z$, the narrow bridge
and long arch to the worker endpoint at $\tau+u$, the reversed worker,
the fixed short turn, the prefix $P$, the giant $I$, and the upper
exterior. Here $z\in[eh/5,eh/3]$ and $u\in[eh,2h/5]$ are relative
heights; $\tau$ is the absolute near-launch level. In the shaded slab
the obstacle is to the left of $x_A-\sigma h$ and the long arch is to
its right. The curves are not to scale and indicate joining order and
clearance, not lattice edges or monotonicity. In particular $I$ has no width bound, and the
drawn prefix endpoint has no required height order relative to $u$.}
\label{fig:10-hairpin-geometry}
\end{figure}


\smallskip\noindent\emph{Representation mass and nearby turn lines.}
For the one-giant construction, let $m(\Gamma)$ count the possible
input tuples and joining choices that produce the rooted output bridge
$\Gamma$. Define
\[
 Z_1=\sum_\Gamma\rho^{|\Gamma|}m(\Gamma),\qquad
 V_1=\sum_{m(\Gamma)>0}\rho^{|\Gamma|},\qquad
 M_2=\sum_\Gamma\rho^{|\Gamma|}m(\Gamma)^2.
\]
The weight of an input tuple is comparable to the output weight,
because its pieces partition $\Gamma$ apart from the bounded local
port changes. Thus $Z_1$ is comparable to the summed construction
weight. Weighted Cauchy--Schwarz gives $Z_1^2\le V_1M_2$.
For $X_d=X_d(h,r)$, the preceding construction and the strip lateral
estimate therefore give
\[
 Z_1\ge cX_d r^{2\alpha}h^{1/2},\qquad
 V_1\le Cr^\alpha\Phi(d).
\]
The factor \(h^{1/2}\) is the product of the \(\asymp h\) choices of
\(z\) and the two \(h^{-1/4}\) connector masses. The marked giant's
width forces comparable output range, giving \(\Phi(d)\).

We estimate the second moment by pairs of representations of the same
output. First bin their turn lines \(\tau\), in absolute output coordinates,
into intervals of width \(r/50\). There are only boundedly many such
bins, because output heights are \(O(r)\). Cauchy--Schwarz permits us to
sum same-bin second moments at bounded cost. For turn-line difference
at most \(Nh\), the contribution is
\begin{equation}
 C_Nr^{2\alpha}X_d(h,r)h^{1/2}h^{5/4}.\label{eq:T-35}
\end{equation}
Here is the full cost allocation. Subdivide into intervals of width
\(ceh\), with a sufficiently small fixed \(c\). For nearby cells,
\(ab\le(a^2+b^2)/2\) reduces pair counting to same-cell pairs with a
factor depending on \(N,e\), but no factor for the absolute cell index.
A line strictly between all \(\tau+z\) and all \(\tau+u\) in the cell has the
following property: after its first visit, the lowest point of the path
is the short turn. Before the descending worker reaches that turn, the
initial rising branch stays above \(\tau+z\); afterwards everything other
than the fixed short arch stays above the turn line. Thus the short
turn, its line, and its two bounding visits coincide in the two
representations. The upward branch determines the giant mark up to
bounded choices, since its prefix is at most \(h/8\). Retain the marked
input associated to the smaller worker height \(u\).

Between the two worker heights there are two bridges, one on either
side of the long arch. Their jointly summed cost is
\(\sum_{j\le Ch}B_j^2\le Ch^{1/2}\). The higher long arch costs
\(Ch^{-1/4}\). The far rising bridge from the higher \(z\) to the smaller
\(u\) has height comparable to \(h\), hence costs \(Ch^{-1/4}\).
The pair of \(z\)-levels and their gap cost
\(Ch\sum_{j\le Ch}B_j\le Ch^{1+\alpha}\). Multiplying these four costs
gives \(Ch^{1/2+5/4}\), as in \eqref{eq:T-35}. The two exterior
continuations supply \(r^{2\alpha}\). Ties use \(B_0=1\).
Every cut is determined by its level: each selected \(u\)-level is
crossed twice on the pre-turn part, while the \(z\)-cuts are first
visits before these crossings. A rooted tuple of these relative shapes
has one translation to the prescribed joining port, also when attached
by its far endpoint. Thus no horizontal multiplicity is omitted.

\smallskip\noindent\emph{Distant turn lines.}
For distant turn lines in the same coarse bin, write
\(\Delta=\tau_2-\tau_1>Nh\), and let \(p_i=\tau_i+f_i\) be their giant junctions.
Everything before \(p_1\) has height at most
\(\tau_1+(2/5+\sigma)h<\tau_2+z_2\), since \(\Delta>Nh\).
Thus the first visit to \(\tau_2+z_2\) occurs after \(p_1\).
Moreover, \(\Delta<r/50\) and \(z_2=O(h)\), with \(h<\Delta/N\),
put this cut below the terminal level of the first giant. The portion
from \(p_1\) to this cut is therefore contained in that giant. It is
strict above its lower level by the giant's lower wall, and strict
below its upper level by the first-visit property. After the cut the
second representation never falls below \(\tau_2-O(1)\). If this subbridge had a
single-crossing level below \(\tau_2-O(1)\), it would also be a seam of the
first giant. Its first irreducible \(I'\) consequently has height
\[
 J(I')\in[\Delta-f_1-O(1),\Delta-f_1+O(h)],
\]
and the rest is a bridge of free height \(O(h)\). Replacing the first
giant by \(I'\) preserves the earlier qualification, which concerns
only its near range. Retain this reduced first construction and then
the second construction with its original giant. The very first and
last exterior continuations are retained as well. The removed endpoint
of the first giant is recoverable as its first ensuing renewal in the
upward continuation; it has no unbounded marking multiplicity.

The second construction costs \(h^{1/2}X_d(h,r)\); the reduced first
one at dyad \(a\) costs \(h^{1/2}X_0(h,a)\); the intervening bridge costs
\(h^\alpha\). Hence the distant contribution is at most
\[
 Cr^{2\alpha}h^{1/2}X_d(h,r)\,
 h^{\alpha+1/2}\sum_{Nh/C\le a<r}X_0(h,a).
\]
Relative offsets of the two turn lines are already determined by the
reduced shapes. For \(d=0\), induct on the dyadic \(r\) with the bound
\(X_0(h,a)\le Kh^\alpha a^{-\alpha}\). After division by
\(r^{2\alpha}h^{1/2}X_0(h,r)h^{5/4}\), the distant contribution is at
most \(CKN^{-\alpha}\). Cauchy--Schwarz, \(Z_1^2\le V_1\) times the
second moment, therefore gives
\[
 X_0(h,r)\le h^\alpha r^{-\alpha}(C_N+CKN^{-\alpha}).
\]
Choose \(N\) first to make the coefficient of \(K\) smaller than one
half, then \(K\) to absorb \(C_N\) and the bounded initial scales. For
\(d>0\) use this proved \(X_0\) bound in the distant sum and retain the
factor \(\Phi(d)\) in \(V_1\). This proves the one-giant bound.

For \(Y\), only nearby bins are needed, with \(r=h\). The obstacle's
terminal cut \(v\) is now movable on the common upward branch. Retain
the smaller worker height and smaller obstacle height, even if they
come from different representations. Shortening preserves
qualification and \(v-u\ge2\sigma h\). The larger \(v\) costs one
additional summed bridge gap, at most \(Ch^\alpha\). The same
Cauchy--Schwarz calculation gives \(Y\le Ch^\alpha\).

\smallskip\noindent\emph{Two ends of a shared giant.}
It remains to bound both ends of one giant simultaneously. Use the
analogous representation count for the paired construction. Its first
moment is at least
\(cX_{12}r^{2\alpha}\prod_i h_i^{1/2}\), and output mass is at most
\(Cr^\alpha\). Bin both turn lines into intervals of width \(r/50\).
On side \(i\), color the fine cells of width \(ceh_i\) periodically with
a sufficiently large fixed number of colors. Matching colors mean either
the same cell or separation exceeding \(100h_i\). Splitting into these
finitely many classes and applying Cauchy--Schwarz lets us use only
matching-color pairs.

Call a same-cell side local. The preceding recovery of the short turn
also recovers its giant junction: all seams within the first \(h_i/8\)
bands are unaffected by the other-end hairpin, which lies far beyond
that range. Distinct junctions would put an internal seam in one of the
two giants. On a distant side the outer junction precedes the inner one
in the inward orientation, and the preceding \(I'\)-replacement applies.

We describe the retained subpaths before counting them. Write
\(a_j,b_j\) for the lower and upper giant junctions of representation
\(j\), and write \(\prec\) for order along the oriented output path.
At the lower end denote the two junctions by
\(a_{\rm out}\preceq a_{\rm in}\), and at the upper end by
\(b_{\rm in}\preceq b_{\rm out}\); equality means that side is local.
The subpath between \(a_{\rm in}\) and \(b_{\rm in}\) will be the
central giant. When both sides are distant, the two possible orders,
after interchanging representation names or reflecting, are
\[
 a_1\prec a_2\prec b_1\prec b_2
 \quad\hbox{and}\quad
 a_1\prec a_2\prec b_2\prec b_1.
\]
Thus the central subpath is \([a_2,b_1]\) in the crossed case and
\([a_2,b_2]\) in the nested case. Figure~\ref{fig:10-paired-path-order}
shows these path orders; it does not depict the spatial positions of
the junctions.
\begin{figure}[htbp]
\centering
\begin{tikzpicture}[x=1cm,y=1cm,font=\small,>=Stealth]
\foreach \offset/\name/\third/\fourth in {0/Crossed/1/2,7.5/Nested/2/1}{
 \begin{scope}[xshift=\offset cm]
  \node[anchor=west] at (0,1.28) {\name};
  \draw[->] (-.25,0)--(6.0,0) node[anchor=west,font=\footnotesize] {path};
  \draw[line width=2.3pt,gray!55] (2,0)--(3.75,0);
  \foreach \x/\lab in {.25/a_1,2/a_2,3.75/b_{\third},5.5/b_{\fourth}}{
   \fill (\x,0) circle (1.4pt);
   \node[anchor=north] at (\x,-.09) {$\lab$};
  }
  \node[anchor=north,font=\footnotesize] at (2.875,-.62) {central giant, counted once};
 \end{scope}
}
\draw[blue!65!black] (.25,.18)--(.25,.75)--(3.75,.75)--(3.75,.18);
\node[fill=white,inner sep=1pt,text=blue!65!black] at (1.0,.75) {$I_1$};
\draw[orange!75!black] (2,.18)--(2,.43)--(5.5,.43)--(5.5,.18);
\node[fill=white,inner sep=1pt,text=orange!75!black] at (4.7,.43) {$I_2$};
\draw[blue!65!black] (7.75,.18)--(7.75,.75)--(13,.75)--(13,.18);
\node[fill=white,inner sep=1pt,text=blue!65!black] at (8.6,.75) {$I_1$};
\draw[orange!75!black] (9.5,.18)--(9.5,.43)--(11.25,.43)--(11.25,.18);
\node[fill=white,inner sep=1pt,text=orange!75!black] at (10.375,.43) {$I_2$};
\end{tikzpicture}
\caption{Two representations of the same output when both sides are
distant, shown in path order, not in spatial coordinates. In the crossed case
$a_1\prec a_2\prec b_1\prec b_2$, the central giant is $[a_2,b_1]$;
in the nested case $a_1\prec a_2\prec b_2\prec b_1$, it is
$[a_2,b_2]$. The brackets identify the original giants $I_j=[a_j,b_j]$
as recovered subpaths. The two outside junctions are retained by the
reduced inputs. Neither the bracket crossings nor their vertical
positions represent intersections or heights of the walk.}
\label{fig:10-paired-path-order}
\end{figure}


The central subpath is indeed an irreducible bridge of scale \(r\) or
\(r/2\). When both its junctions come from one representation, it is
that representation's giant. In the crossed case it is contained in
both giants: the giant starting at \(a_2\) keeps it strictly above
\(a_2\), and the giant ending at \(b_1\) keeps it strictly below
\(b_1\). Divide its height range into lower and upper halves.
A seam in the lower half would extend to a seam of the
giant \([a_2,b_2]\): between \(b_1\) and \(b_2\), the upper
hairpin can descend below \(b_1\) by at most
\((2/5+\sigma)h_{\rm up}+O(1)\), where \(h_{\rm up}\) is the
upper-side test scale. A seam in its upper half is
excluded by the reflected argument using \([a_1,b_1]\). Each turn-line
difference is less than \(r/50\), and a distant difference exceeds
\(100h_i\), so the near-end excursions cannot invalidate either
extension. The central height can fall below \(r\) only in this
two-distant case, where necessarily \(h_i\le r/2\). On a distant
side retain the inner representation's worker and prefix; on a local
side retain the smaller worker height. They qualify against the central
giant by containment, and so give a central paired input.

\smallskip\noindent\emph{Recovering the paired representations.}
In the crossed case, the lower reduced construction retains the first
representation's worker and prefix, followed by its first irreducible \(I'_-\); thus
\(a_1\) is still the prefix--irreducible junction of that counted
input. On the upper side, read inward from the upper end, the reduced
construction retains the second representation's worker and prefix,
followed by its first irreducible \(I'_+\); its prefix--irreducible
junction is \(b_2\). The remaining endpoints \(a_2,b_1\) are the two
endpoints of the central giant. These are existing input junctions,
rather than additional marks placed on an unmarked bridge.

Here is the order in which the lower pieces are joined. For
\(j=1,2\), write \(A_j,P_j,S_j\) for the lower worker, obstacle
prefix, and fixed short turn, respectively. Let \(K_j\) be the
connector from the cut at \(\tau_j+z_j\), through the narrow
bridge and long arch, to the worker endpoint, retaining their counted
joining ports and local connection data; here
\(\tau_j\) denotes the turn-line height. Let \(R_-\) be the
bridge remainder after \(I'_-\), and let \(E_-\) be the first
exterior continuation. With a bar denoting path reversal, the lower
part of the common output is assembled in the order
\[
 E_-\,K_1\,\overline{A_1}\,S_1\,P_1\,I'_-\,R_-\,
 K_2\,\overline{A_2}\,S_2\,P_2,
\]
and then the central giant is attached. The upper part is assembled
by the reflected reverse procedure, with representation 2 outer and
representation 1 inner. In the forward cutting operation,
\(I'_-\) ends at the first positive renewal of the strict subbridge
from \(a_1\) to the second connector's first-visit cut; \(R_-\)
is the rest of that subbridge. The upper cut uses the same rule in
the reversed orientation. Consequently the reduced-irreducible cuts
require no further renewal index. The connector cuts themselves are
the first-visit and ordered worker cuts already specified above.

Each item in this list is a rooted relative shape: part of the central
or a reduced outer input, a connector, a bridge remainder, or an
exterior continuation. Once one joining port is placed, the next shape
has a unique translation to it, including when the shape is attached
by its terminal port. Its height and horizontal displacement are then
fixed. Starting at the prescribed output root therefore places every
piece, and in particular all four retained junctions, without a
separate sum over turn-line differences or horizontal offsets.
The second lower exterior continuation is the assembled prefix
\(E_-K_1\overline{A_1}S_1P_1I'_-R_-\) ending at the exterior
port of \(K_2\);
the remaining upper exterior is the corresponding recovered suffix.
Thus each original exterior is recovered by cutting at its retained
connector port. Finally, extract the ordered subpaths \([a_1,b_1]\) and
\([a_2,b_2]\), and assign the retained workers, prefixes and
connectors to their recorded representation. This recovers both
original constructions. Interchanging the representation names,
choosing the reflected orientation, and the bounded local connection
choices require only finitely many records.

For the nested two-distant case, the order is
\(a_1\prec a_2\prec b_2\prec b_1\). The central paired input is
bounded by \(a_2,b_2\), both from representation 2, and the two
reduced inputs retain \(a_1,b_1\), both from representation 1.
A finite matching record assigns the four retained boundaries
\(a_{\rm out},a_{\rm in},b_{\rm in},b_{\rm out}\)
to the two representations; the same ordered assembly and exterior-port cuts recover the nested
case. On a local side the two junctions coincide as proved above.
The tuple pieces partition the common output, apart from bounded
port modifications. Their product critical weight is therefore
comparable to the single critical weight of that output in the
second moment. In particular the central giant is charged once;
the two original giants are recovered subpaths of the assembled
output and do not contribute additional weight factors.

The reconstruction has bounded ambiguity, uniformly in all scales.
We can therefore integrate these pieces. Each exterior continuation
costs \(Cr^\alpha\). A local side has the four gap and arch costs
already computed, namely \(Ch_i^{1/2+5/4}\). On a distant side the
inner connector costs \(Ch_i^{1/2}\), the reduced outer construction
costs \(Ch_i^{1/2}X_0(h_i,a)\), and the bridge remainder costs
\(Ch_i^\alpha\). The proved one-sided estimate gives
\[
 h_i^{1/2}h_i^{1/2}h_i^\alpha
       \sum_{a\ge c h_i}X_0(h_i,a)
 \le Ch_i^{1+\alpha}=Ch_i^{1/2+5/4}.
\]
The central paired input has scale \(r\) or \(r/2\), and its giant
weight is charged once. Combining these costs with the first moment
and output mass yields
\[
 \begin{split}
 &\left(cX_{12}(h_1,h_2,r)r^{2\alpha}\prod_i h_i^{1/2}\right)^2\\
 &\qquad\le Cr^{3\alpha}\prod_i h_i^{1/2+5/4}
  \{X_{12}(h_1,h_2,r)+X_{12}(h_1,h_2,r/2)\},
 \end{split}
\]
omitting an out-of-domain last term. Writing
\(V_r=X_{12}(h_1,h_2,r)/(h_1^\alpha h_2^\alpha r^{-\alpha})\),
this becomes
\(V_r^2\le C(V_r+2^\alpha V_{r/2})\). A sufficiently large constant
majorant closes the dyadic induction, proving the paired bound.
\end{proof}

\subsection{Finding separated renewal endpoints}

\begin{lemma}[Usable renewals after a regular history]
\label{lem:10-usable-renewals}
Stop a worker at its first renewal at or above \(s\), where \(s/h\) is
a sufficiently small fixed constant. Call this history regular when
the sum of \(J+W\), including its last jump, is less than \(h/16\).
For every \(\delta>0\) one can choose fixed qualification parameters so
that, for all sufficiently large \(h\) depending on these choices,
conditionally on any such history and any fixed opposite path,
with failure probability at most \(\delta\), continued \(E(h)\) implies
at least \(c h^\alpha\) qualified worker renewals in \([eh,2h/5]\).
For two workers with regular histories, outside total conditional failure
probability at most \(2\delta\), continued \(E(h)\) supplies at least
\(ch^{2\alpha}\) separated renewal pairs. In each pair the lower renewal
is a qualified worker against the opposite prefix ending at the higher
renewal, whose height exceeds it by at least \(2\sigma h\).
The failure probabilities are taken before conditioning on continued
\(E(h)\).
\end{lemma}

\begin{proof}
Write \((t,x_0)\) for the worker stop. A regular history has
\(t<h/16\) and horizontal range less than \(h/16\). For
\(m=\lceil3c'h^\alpha\rceil\), the probability that the next \(m\)
jumps advance more than \(h/4\) is \(O(c')\): bound jumps larger than
\(h/4\) by a union bound and the remaining sum by its truncated mean.
Choose \(c'\) small. There are then at least \(2c'h^\alpha\) future
renewals before \(t+h/4<2h/5\), except with arbitrarily small
probability. Renewals within \(eh\) of the stop have expected count
\(O((1+eh)^\alpha)\), so remove them by taking \(e\) small.

No short bin contains a substantial fraction of these renewals with
high probability. Indeed partition the candidate interval into
\(O(\theta^{-1})\) bins of length \(\theta h\). By
Lemma~\ref{lem:10-charges}, for any fixed \(p>1/\alpha\), the probability
that some bin contains at least \(\varepsilon h^\alpha\) renewals is at
most \(C_{p,\varepsilon}\theta^{p\alpha-1}\), once \(\theta h\) is
large. Use neighboring bins as well to control every interval of length
\(2\theta h\). The charge bound also makes the entire worker's lateral
range through its test stop at \(h\) at most \(Dh\), with failure
probability \(O(D^{-\alpha})\).

Freeze the opposite path. For a right worker let \(R(u)\) be the
rightmost opposite abscissa in \([u,u+\sigma h]\). From \((t,x_0)\)
the renewal-occurrence density at \((u,x)\) is exactly
\(b_{u-t}(x-x_0)\). It is not divided by \(B_{u-t}\), since we count
renewal occurrences. Thus the expected number of renewals, at distance
at least \(eh\) above the stop, with \(|x-R(u)|\le2e'h\), is at most
\begin{align*}
 C(e'h+1)\sum_{j\ge eh}\|b_j\|_\infty
 &\le C(e'h+1)\sum_{j\ge eh}(B_j-B_{j+1})\\
 &\le C_e(e'+h^{-1})h^\alpha.
\end{align*}
The middle inequality follows from Lemma~\ref{lem:08-variation}: a summable nonnegative kernel tends to zero at infinity, so its supremum is at most its total variation. The sum of the resulting differences telescopes. No height-by-height derivative bound on $B_j$ is needed.
At each renewal \(u\), restart and discard it if the path through its
first renewal strictly above \(u+\sigma h\) moves more than \(e'h\)
laterally. Its conditional probability is at most
\(C((1+\sigma h)/(e'h))^\alpha\); multiplying by the expected candidate
count bounds the expected number discarded by this quantity times
\(Ch^\alpha\).

On continued order, every opposite point in this slab lies to the left
of the worker's rightmost point at that level. For a retained candidate
that point is at most \(x+e'h\), because strict renewals exclude earlier
and later pieces from the slab. Hence \(R(u)\le x+e'h\). Since
\(|x-R(u)|>2e'h\), necessarily \(R(u)<x-2e'h\). Choose
\(\sigma<e'\). This is the required clearance; the obstacle's global
one-sided bound below \(u+\sigma h\) follows from order and the worker's
bounded range. Reflection gives the other orientation.

Choose the constants successively: the desired failure tolerance,
\(c'\), then \(e,\theta\), then \(D\) large and \(e'\) small, and finally
\(\sigma\) small relative to \(e',e,\theta\). The expectation bounds and
Markov's inequality leave \(ch^\alpha\) candidates. For two workers,
apply this uniform estimate first with the opposite future frozen and
then with roles reversed, integrating by Tonelli. A union bound suffices;
independence of the two good events is unnecessary. Bin nonconcentration
ensures that only a small fraction of all cross-pairs have heights
within \(2\sigma h\). The remaining pairs have the claimed lower/higher
qualification.
\end{proof}

The preceding lemma converts a regular stopped history into many
qualified endpoints. To use it without biasing the fresh jumps, we
will retain a smaller test that is measurable at the stopping time,
then bound failures under the original conditional law. Before doing
so, we record the dyadic estimates that control irregular histories.

\subsection{The ordered-probe induction}

We record the elementary dyadic sums that close the proof. Let \(a\)
range over positive dyadic heights and let \(l\) range over
\(0,2,4,8,\ldots\); the bin \(l=0\) means \(W/a<2\). Put
\[
 T_h(a,l)=\min\{1,(1+l)a/h\}.
\]
Since \(\sum_l(1+l)\Phi(l)<\infty\) on these dyads,
\begin{align}
 \sum_{a,l}a^{-\alpha}\Phi(l)T_h(a,l)&\le Ch^{-\alpha},
 \label{eq:10-charge-sum}\\
 \sum_{a<\kappa h,l}a^{-\alpha}\Phi(l)T_h(a,l)
        [1+\log^+(\kappa h/a)]
 &\le C\kappa^{1-\alpha}h^{-\alpha}.
 \label{eq:10-small-charge-sum}
\end{align}
For example, first sum in \(l\), obtaining
\(C\min(1,a/h)\), then sum the geometric series below and above \(h\).
The logarithm in the second line is summable against
\(a^{1-\alpha}\). Also, if \(s=h/L\),
\begin{equation}
 (1+s)^\alpha\sum_{a,l}a^{-\alpha}\Phi(l)T_h(a,l)
 [\log(2L)+\log^+((1+s)/a)]
 \le CL^{-\alpha}\log(2L)
 \label{eq:10-comparable-charge-sum}
\end{equation}
for \(h\) sufficiently large relative to the fixed \(L\). The extra
logarithm for \(a<s\) contributes \(O(s/h)\), which is bounded by the
right-hand side.

\begin{proof}[Proof of Theorem~\ref{thm:10-ordered-probes}]
We first prove the \(E\) and \(F\) bounds simultaneously by strong
induction on their test height, with \(r,d\) unrestricted. Floors below
one mean the trivial test, with constant cost. All uses of nontrivial
smaller tests have strictly smaller test height.

For \(F(h,r,d)\), stop the worker at height \(s=h/L\). If its history
is irregular, then
\[
 \mathbf1_{\rm irregular}
 \le16\sum_{k:U_{k-1}<s}\min\{1,(J_k+W_k)/h\}.
\]
Integrate this inequality in \eqref{eq:10-F-integral}, retaining the
large test only where it is needed to imply a smaller test. Let \(f,g\)
be the dyads of one plus the birth depths of the main giant and a
charged worker jump; let \(a,l\) be that worker jump's height and width
bins. Its charge is at most \(CT_h(a,l)\). Both birth dyads are at most
\(C(1+s)\). Constants in what follows do not grow with \(L\).

When \(f<g/(16L)\), retain the whole smaller event
\(F(\lfloor g/2\rfloor,r,d)\). Do not also sum separately over \(f\).
On contributing paths the worker's stop for this smaller test is no
later than the charged birth, which itself is a renewal. Restarting
there, the expected total charges born in the \(g\)-bin are at most
\(Cg^\alpha h^{-\alpha}\), by Lemma~\ref{lem:10-charges}. After dropping
the large test, this is a uniform bound on the conditional future
integral. Hence the contribution, divided by
\(r^{-\alpha}\Phi(d)\), is at most
\[
 CK_2h^{-\alpha}\sum_{g\le C(1+s)}g^\alpha
 \le CK_2L^{-\alpha}.
\]

When \(g<\min(f,a)/(16L)\), retain instead
\(F(\lfloor\min(f,a)/2\rfloor,a,l)\), with the workers' roles exchanged.
The charged jump is now the unique marked jump of that smaller test;
there is no separate \(g\)-birth factor. The original giant is born
later than its arm's smaller test stop. Its subsequent birth count in
the \(f\)-bin is at most \(Cf^\alpha\), uniformly in the stop, including
a possible birth at the stop itself. The ensuing fresh jump costs
\(Cr^{-\alpha}\Phi(d)\). Thus the relative contribution is bounded by
\[
 CK_2\sum_{f\le C(1+s)}f^\alpha
      \sum_{a,l}a^{-\alpha}\Phi(l)T_h(a,l)
 \le CK_2L^{-\alpha}.
\]
This formula is a product of a smaller tested integral, a renewal count,
and a fresh-jump integral; it is not a conditional probability estimate
under the larger \(F\) event.

In the remaining case retain
\(E(\lfloor\min(f,g)/2\rfloor)\). Both subsequent births are renewals
at or beyond their respective smaller stops. Once their earlier
histories are fixed, their count integrals are bounded by
\(Cf^\alpha\) and \(Cg^\alpha\), and the new jump integrals are
independent. The smaller test costs
\(CK_1\min(f,g)^{-\alpha}\), leaving \(CK_1\max(f,g)^\alpha\).
The restrictions defining this case are
\(f\ge g/(16L)\) and \(g\ge\min(f,a)/(16L)\).
Their dyadic sum obeys
\[
 \sum_{f,g}\max(f,g)^\alpha
 \le C(1+s)^\alpha
       [\log(2L)+\log^+((1+s)/a)].
\]
Indeed for \(f\le a\) the two birth dyads are within factor \(16L\),
so at each larger dyad there are \(O(\log(2L))\) partners. For
\(f>a\), the permitted smaller partner is at least \(a/(16L)\),
adding at most \(O(\log^+((1+s)/a))\) choices. Geometric summation of
the larger dyad proves the displayed bound. Equation
\eqref{eq:10-comparable-charge-sum} makes the relative contribution at
most \(CK_1L^{-\alpha}\log(2L)\).

Consider regular histories next. To apply
Lemma~\ref{lem:10-usable-renewals} without paying the unrestricted prefix
mass, retain a smaller-test envelope. If the main birth is before
\(s/(4L)\), use \(F(\lfloor s/2\rfloor,r,d)\). Otherwise use
\(E(\lfloor s/(8L)\rfloor)\), followed by the main birth count
\(O(s^\alpha)\) and its jump weight. In the later-birth branch the obstacle's smaller
stop occurs no later than the counted birth. In both branches the
worker's smaller stop occurs no later than its stop at \(s\). The total envelope mass is therefore at most
\[
 C_L(K_1+K_2)r^{-\alpha}\Phi(d).
\]
Condition on the worker history to \(s\) and on the full opposite
obstacle before invoking the lemma. Its failure probability is uniformly
at most \(\delta\), without conditioning on continued \(E(h)\).
The envelope is measurable with respect to this exposed information;
integrating it bounds the failure contribution by \(\delta\) times the
last display.

On success there are at least \(ch^\alpha\) qualified renewal choices.
Count them with \(\mathcal R\), using its occurrence identity. Each
choice leaves a qualified worker bridge and the original marked
obstacle; all unused worker futures have probability at most one.
Tonelli therefore bounds the total choice mass by \(X_d(h,r)\).
Lemma~\ref{lem:10-hairpins}, divided by the guaranteed number of choices,
bounds this success contribution by \(Cr^{-\alpha}\Phi(d)\). Altogether,
\begin{equation}
 \frac{F(h,r,d)}{r^{-\alpha}\Phi(d)}
 \le C+CL^{-\alpha}[K_2+K_1\log(2L)]
        +\delta C_L(K_1+K_2).
 \label{eq:10-F-recursion}
\end{equation}

For \(E(h)\), stop both workers at \(s'=h/L'\), where \(L'\) will be
large compared to \(L/\kappa\). If either early history has a jump of
height dyad \(a\ge\kappa h\), retain the smaller event
\(F(\lfloor\kappa h\rfloor,a,0)\); its birth is early enough since
\(s'<\kappa h/L\). Summing over \(a\ge c\kappa h\) costs at most
\(CK_2\kappa^{-\alpha}h^{-\alpha}\).
For all other charged jumps, \(a<\kappa h\). If their birth is before
\(a/L\), use \(F(a,a,l)\), with no further birth-count factor.
Otherwise retain a smaller \(E\)-test below the birth dyad and then
count subsequent births and the new jump as above. The number of
possible late birth dyads contributes
\(1+\log^+(s'L/a)\). Consequently the irregular-history contribution,
relative to \(h^{-\alpha}\), is at most
\begin{align*}
 Ch^\alpha\sum_{a<\kappa h,l}
 &[K_2+K_1\{1+\log^+(s'L/a)\}]
          a^{-\alpha}\Phi(l)T_h(a,l)\\
 &\le C(K_1+K_2)\kappa^{1-\alpha},
\end{align*}
by \eqref{eq:10-small-charge-sum} and \(s'L\le\kappa h\).

For regular histories retain the envelope \(E(\lfloor s'\rfloor)\),
whose mass is at most \(CK_1(s')^{-\alpha}\). Apply the usable-renewal
lemma to each worker while freezing the other path; integrate each
failure estimate and take a union bound. Their total cost is
\(C\delta K_1(s')^{-\alpha}\). On success count the
\(ch^{2\alpha}\) separated cross-pairs. Each pair, according to which
renewal is lower, contributes to \(Y(h)\) in one of two reflected
orientations. The occurrence identity integrates unused futures to at
most one, so the total success probability is at most
\(CY(h)/h^{2\alpha}\le Ch^{-\alpha}\). Thus
\begin{equation}
 \frac{\mathbf P(E(h))}{h^{-\alpha}}
 \le C+CK_2\kappa^{-\alpha}
       +C(K_1+K_2)\kappa^{1-\alpha}
       +C\delta K_1(L')^\alpha.
 \label{eq:10-E-recursion}
\end{equation}

Choose \(\kappa\) so small that the coefficient of
\((K_1+K_2)\) in the third term of \eqref{eq:10-E-recursion} is small.
Choose the ratio \(K_1/K_2\) sufficiently large compared to
\(\kappa^{-\alpha}\). Next choose \(L\) so that the charge terms of
\eqref{eq:10-F-recursion} are a small fraction of \(K_2\), and then
\(L'\gg L/\kappa\). Choose the two failure tolerances small enough
for these fixed constants. Finally increase \(K_1,K_2\) together in
their fixed ratio to absorb the success constants and bounded test
heights. At bounded heights, the prefix sum is bounded and the giant
tail supplies \(r^{-\alpha}\Phi(d)\). Both recursions close.

It remains to prove the paired bound. Regard the two established
constants as fixed and induct on \(h_1+h_2\), using
\(s_i=h_i/L_2\). If either test height is bounded, integrate its prefix
and worker at bounded cost and use the single \(F\) bound on the other
side; since \(L_2\ge L\), its prefix restriction is stronger than that
of \(F\). This gives the initial cases.

For an irregular worker on side 1, use the same three birth-order
regimes, now with \(L_2\). In the very early main-birth regime, retain
the smaller paired test on side 1 and the original side-2 test. Restart
only the side-1 worker after its smaller stop. This gives relative charge
\(CK_3L_2^{-\alpha}\). In the very early worker-birth regime, the
smaller single \(F\) marks the charged worker jump and only uses side
1's obstacle before its later shared-giant birth. Once this tested
prefix and the subsequent side-1 birth count have been integrated, the
remaining shared-giant and side-2 integral is
\[
 \sum_{r\le J(I)<2r}\nu(I)
 \int_{f_2<h_2/L_2}\!d\mathcal R(P_2)
 \int\mathbf1_{E(h_2)}(P_2I^{(2)},A_2)\,d\mathcal Q_{h_2}(A_2)
 \le K_2r^{-\alpha}.
\]
It is bounded by the single \(F\) mass because \(L_2\ge L\). The
pre-giant side-1 factors are independent of this remaining integral:
translating \(I\) to the endpoint of the side-1 prefix adds no weight
and no choice. This yields relative charge \(CK_2^2L_2^{-\alpha}\).
In the comparable-birth regime the smaller side-1 test is \(E\), followed
by its two birth counts and fresh charged jump, then exactly the same
remaining integral. Its charge is
\(CK_1K_2L_2^{-\alpha}\log(2L_2)\).
The giant has been charged once in all three regimes. The full side
configurations are not independent; their only dependence left after
peeling the pre-giant factors is precisely the one shared giant in the
last integral. Repeating the argument on side 2 gives total relative
bad-history cost
\[
 CL_2^{-\alpha}
       [K_3+(K_2^2+K_1K_2)\log(2L_2)].
\]

For regular histories keep measurable smaller-test envelopes before
integrating failure probabilities. If both main births are very early,
retain a paired test with heights \(\lfloor s_i/2\rfloor\). On each
later-birth side retain instead
\(E(\lfloor s_i/(8L_2)\rfloor)\) and its subsequent prefix count.
If one early side remains, keep its smaller single \(F\) test. If
neither remains, charge the giant by its one-jump tail. Thus the total
exposed envelope mass is at most
\[
 C_{L_2}(K_3+K_1K_2+K_1^2)r^{-\alpha}.
\]
All retained worker test stops occur no later than the corresponding
worker stops at \(s_i\); retained obstacle data belong to the exposed
paths. Hence these envelopes are measurable before either tested
worker future is integrated. Apply the conditional failure bound
separately to the two independent worker futures and integrate by
Tonelli; a union bound makes its cost at most \(\delta\) times a fixed
multiple of this display. This does not assert independence of the
usability events or condition either future on continued survival.

On success each side has at least \(ch_i^\alpha\) usable renewals.
Their simultaneous choices number at least
\(c h_1^\alpha h_2^\alpha\), since there is no condition between the
two side tests. In \eqref{eq:10-paired-integral}, sum both renewal
occurrences before integrating. The shared giant remains a single
factor \(\nu(I)\); each worker prefix becomes its critical bridge
measure and its unused future has mass at most one. The resulting
integral is bounded by \(X_{12}(h_1,h_2,r)\), so the paired hairpin
bound makes the success contribution at most \(Cr^{-\alpha}\).
We have proved
\[
 \frac{F_{12}(h_1,h_2,r)}{r^{-\alpha}}
 \le C+CL_2^{-\alpha}
       [K_3+(K_2^2+K_1K_2)\log(2L_2)]
       +\delta C_{L_2}(K_3+K_1K_2+K_1^2).
\]
Choose \(L_2\ge L\) large, then \(\delta\) small, and finally \(K_3\)
large. This closes the induction and proves \eqref{eq:T-34}.
\end{proof}

\subsection{Avoidance up to the first hit}

The next application forgets the order condition in the definition of
$E(h)$ and stops each path when it first visits a level, rather than
at a renewal above that level. Define $S_2(h)$ as the probability that
two independent infinite renewal walks from neighboring horizontal
ports are disjoint up to their respective first visits to height $h$.
The renewal stops may lie beyond this height. We handle those
overshoots explicitly.

\Needspace{5\baselineskip}
\begin{corollary}[First-hit avoidance with logarithmic loss]
\label{cor:10-log-first-hit}
For every integer $h\ge1$, the first-hit probability just defined satisfies
\begin{equation}
 S_2(h)\le Ch^{-\alpha}(1+\log h)^2.\label{eq:T-35a}
\end{equation}
\end{corollary}

\begin{proof}
Stop both renewal strings at their first renewals at or above \(h/4\).
If both stops are below \(h\), disjoint first-hit arcs to \(h\) imply
\(E(\lfloor h/4\rfloor)\). Indeed those renewal stops are single
crossings, so later portions cannot revisit the tested lower region;
the two first-hit crosscuts have their launch order throughout it.
Equation \eqref{eq:T-34} bounds this case.

Otherwise mark each jump born below \(h/4\) that reaches or exceeds
\(h\). There are one or two such marked jumps, each of height at least
\(3h/4\). Bin their birth depths plus one into dyads \(f_i\). Retain
the ordered test at half the smallest birth dyad, or the trivial test
at bounded depth. Its stops precede the corresponding births, and in
the unmarked arm precede its stop below \(h\), so they precede both
first hits at \(h\). The same crosscut argument makes this test valid.
It costs \(Cf_{\min}^{-\alpha}\). Each subsequent marked birth count
costs \(Cf_i^\alpha\) by restarting, and each ensuing jump costs
\(Ch^{-\alpha}\) by the irreducible tail. Thus one marked bin costs
at most \(Ch^{-\alpha}\), and a pair of marked bins costs at most
\[
 C f_{\min}^{-\alpha}(f_1/h)^\alpha(f_2/h)^\alpha
 =C f_{\max}^\alpha h^{-2\alpha}\le Ch^{-\alpha}.
\]
There are \(O(1+\log h)\) bins per marked jump. Summing proves
\eqref{eq:T-35a}.
\end{proof}

\begin{corollary}[Adjacent-arm avoidance]\label{09_fast_paths:adjacent-arms}
For every \(\epsilon>0\), the probability \(S_2(H)\) of disjoint
first-hit paths from neighboring ports satisfies
\(S_2(H)\le C_\epsilon H^{-3/4+\epsilon}\).
\end{corollary}
\begin{proof}
Corollary~\ref{cor:10-log-first-hit} gives
\(S_2(H)\le CH^{-3/4}(1+\log H)^2\). For fixed \(\epsilon>0\),
the logarithmic factor is bounded by \(C_\epsilon H^\epsilon\).
\end{proof}
The weaker power estimate is enough for the fixed number of turns in the
next section. Appendix~\ref{sec:asynchronous} gives its independent proof
by extending asynchronous renewal pairs to a common height.
